
\section{Diseño del enfoque}
\label{sec:diseno}

Antes de entrar en los detalles, conviene fijar la lógica que guía el
enfoque. En esta sección se presenta la estrategia general y los
principios que la sostienen; a continuación, el patrón objetivo y las
precondiciones que delimitan cuándo es seguro aplicarlo; y, por último, la
política de transformación y la política iterativa con las que el
prototipo opera sobre el código.

\subsection{Estrategia general}
El enfoque responde a \textbf{RQ1} mediante un prototipo determinista que
trabaja directamente sobre el árbol de sintaxis abstracta (AST) de Java.
La decisión de partida es deliberadamente \textbf{conservadora}: ante la
menor duda sobre la equivalencia semántica de una transformación, el caso
se rechaza. Es preferible perder oportunidades de refactorización antes que
arriesgar un cambio de comportamiento, porque el prototipo aspira a servir
como baseline fiable frente al cual medir a los LLMs en RQ3.

Esta estrategia se concreta en tres principios:

\begin{enumerate}
  \item \textbf{Aplicabilidad explícita.} Toda transformación queda
  condicionada a un conjunto cerrado de precondiciones formales, evaluadas
  sobre el AST. Si alguna no se cumple, el caso se descarta con un motivo
  trazable, lo que permite justificar después por qué un fragmento no se
  combinó.
  \item \textbf{Preservación del orden de evaluación.} La condición
  combinada respeta el orden original \lstinline|A && B|. No se reordenan
  los operandos ni se simplifica el álgebra booleana, de modo que la
  evaluación en cortocircuito se comporta igual que antes.
  \item \textbf{Determinismo y reproducibilidad.} Para una misma entrada,
  detector y transformador producen siempre la misma salida. Es un
  requisito imprescindible para que el contraste con los LLMs en RQ3 sea
  justo y repetible.
\end{enumerate}

En conjunto, estos principios trazan una frontera nítida entre lo que el
prototipo se atreve a transformar y lo que prefiere dejar intacto, frontera
que las precondiciones de la \autoref{sec:precondiciones} formalizan.

\subsection{Patrón objetivo de transformación}
El patrón canónico (\autoref{fig:before-after}) es semánticamente equivalente
al original \emph{bajo} las precondiciones estructurales P1--P4. La equivalencia
se sostiene en el cortocircuito de \lstinline|&&|~\cite{jls17}: como \lstinline|A && B|
equivale a \lstinline|A ? B : false|, evalúa \lstinline|A| y, solo si es cierto,
\lstinline|B|, exactamente igual que el anidamiento. Por ello la combinación
preserva el comportamiento observable \emph{con independencia de que las
condiciones tengan efectos colaterales}: estos ocurren en el mismo orden y bajo
las mismas condiciones en ambas versiones. Las precondiciones que realmente
importan son, por tanto, las puramente estructurales (P1--P4); no se requiere
análisis de pureza ni de flujo de datos para garantizar la corrección.

\subsection{Precondiciones de aplicabilidad (P1--P4)}
\label{sec:precondiciones}
El detector evalúa, para cada nodo \lstinline|IfStmt| del AST, cuatro
precondiciones \emph{estructurales}, necesarias y conjuntamente suficientes
para considerar el caso elegible (la razón de la suficiencia, el cortocircuito
de \lstinline|&&|, se explica en la \autoref{sec:patron}). Los nombres en monoespaciado como
\lstinline|IfStmt| o \lstinline|BlockStmt| son
tipos de nodo del AST de JavaParser, no clases del prototipo; las clases y
enumeraciones propias del prototipo (incluidos los motivos de descarte) se
recogen en el diagrama de clases (\autoref{fig:clases-nucleo}).

\begin{description}
  \item[P1 — Ausencia de \texttt{else} externo.] El \lstinline|if| externo
  no tiene rama \lstinline|else| ni \lstinline|else if|. Un flujo
  alternativo no se preserva con la combinación \lstinline|A && B|.
  \item[P2 — Rama externa unitaria.] La rama \lstinline|then| del
  \lstinline|if| externo consiste en \emph{exactamente} una sentencia, ya
  sea encerrada en un bloque \lstinline|{...}| de una sola sentencia o
  escrita directamente sin llaves.
  \item[P3 — La sentencia única es un \texttt{if}.] La única sentencia de la
  rama \lstinline|then| exigida por P2 debe ser, a su vez, otro
  \lstinline|IfStmt| (el \lstinline|if| interno candidato). El detector
  contempla ambas formas, con y sin llaves (p.\,ej.\
  \lstinline|if (A) if (B) S;|).
  \item[P4 — Ausencia de \texttt{else} interno.] El \lstinline|if| interno
  tampoco tiene rama \lstinline|else| ni \lstinline|else if|.
\end{description}

En conjunto, P1 y P4 exigen que ninguno de los dos \lstinline|if| (ni el
externo ni el interno) tenga rama \lstinline|else| o \lstinline|else if| ---que
en el AST de JavaParser no es una construcción propia, sino un \lstinline|else|
cuyo cuerpo es a su vez un \lstinline|IfStmt|, de modo que la comprobación de
ausencia de \lstinline|else| (\lstinline|hasElseBranch()|) cubre ambos casos por
igual---, y P2 y P3, que el bloque externo contenga exactamente el
\lstinline|if| interno y nada más. Estas cuatro condiciones son puramente estructurales y, por el
cortocircuito de \lstinline|&&| (\autoref{sec:patron}), necesarias y
suficientes para que la combinación preserve el comportamiento.

La \autoref{fig:arbol-precondiciones} resume el flujo de decisión completo: las
precondiciones se evalúan en orden y, en cuanto una falla, el candidato se
descarta con un motivo trazable; solo si todas se cumplen el caso se declara
elegible.

\begin{figure}[ht]
\centering
\begin{tikzpicture}[
  check/.style={draw, rounded corners, fill=blue!5, align=center,
                text width=4.2cm, inner sep=4pt, font=\small},
  reject/.style={draw, fill=red!5, align=left,
                 text width=5.0cm, inner sep=3pt, font=\scriptsize},
  accept/.style={draw, thick, fill=green!12, align=center,
                 text width=4.2cm, inner sep=4pt, font=\small},
  arr/.style={-{Latex}},
  lbl/.style={font=\scriptsize, inner sep=1pt},
]
% Columna izquierda (decisiones), coordenadas explícitas
\node[check]  (p1) at (0,0)    {P1 — \textquestiondown{}\texttt{if} externo sin \texttt{else}?};
\node[check]  (p2) at (0,-2)   {P2 — \textquestiondown{}rama \texttt{then} con una sola sentencia?};
\node[check]  (p3) at (0,-4)   {P3 — \textquestiondown{}esa sentencia es un \texttt{if}?};
\node[check]  (p4) at (0,-6)   {P4 — \textquestiondown{}\texttt{if} interno sin \texttt{else}?};
\node[accept] (ok) at (0,-8)   {\textbf{Elegible:} combinar \texttt{A \&\& B}};
% Columna derecha (motivos de descarte)
\node[reject] (r1) at (8,0)    {\texttt{OUTER\_HAS\_ELSE}};
\node[reject] (r2) at (8,-2)   {\texttt{OUTER\_BLOCK\_MULTIPLE\_STATEMENTS}};
\node[reject] (r3) at (8,-4)   {\texttt{SINGLE\_STATEMENT\_NOT\_IF}};
\node[reject] (r4) at (8,-6)   {\texttt{INNER\_HAS\_ELSE}};
% Flujo "sí" (vertical)
\draw[arr] (p1) -- node[lbl,left]{sí} (p2);
\draw[arr] (p2) -- node[lbl,left]{sí} (p3);
\draw[arr] (p3) -- node[lbl,left]{sí} (p4);
\draw[arr] (p4) -- node[lbl,left]{sí} (ok);
% Flujo "no" (horizontal hacia el motivo de descarte)
\draw[arr] (p1) -- node[lbl,above]{no} (r1);
\draw[arr] (p2) -- node[lbl,above]{no} (r2);
\draw[arr] (p3) -- node[lbl,above]{no} (r3);
\draw[arr] (p4) -- node[lbl,above]{no} (r4);
\end{tikzpicture}
\caption{Árbol de decisión del detector. Las precondiciones estructurales
P1--P4 son necesarias y suficientes para la combinación; cada rama "no" produce
un motivo de descarte estable (\texttt{DiscardReason}).}
\label{fig:arbol-precondiciones}
\end{figure}

Estas precondiciones están implementadas en \texttt{NestedIfDetector} y
materializadas como categorías estables del enumerado
\texttt{DiscardReason}, que el detector emite al rechazar un candidato. El
catálogo completo de motivos (incluyendo la categoría auxiliar
\texttt{NO\_NESTED\_IF\_PATTERN}) se reproduce en el
Apéndice~D. El detector expone dos métodos de salida:
\texttt{detect(method)} (solo oportunidades aceptadas) y
\texttt{detectWithReasons(method)} (también los candidatos rechazados con
sus motivos), siendo este último el que alimenta la columna
\texttt{discardCategories} del CSV de resultados y da respuesta empírica a
RQ1.

El detector aplica estas cuatro precondiciones en el modo \texttt{STRUCTURAL}
del enumerado \texttt{DetectionMode}, que es el que sustenta los
\textbf{resultados principales} de esta memoria.

\subsection{Política de transformación y seguridad semántica}

Una vez detectada una oportunidad, \texttt{NestedIfTransformer} aplica la transformación sobre el AST siguiendo esta secuencia:

\begin{enumerate}
    \item \textbf{Guarda previa.}
    Al entrar, el transformador comprueba que ninguno de los dos \lstinline|if| tenga rama \lstinline|else|; si la tuviera, aborta devolviendo \lstinline|false| sin modificar el AST. Esta comprobación es redundante respecto a P1 y P4, pero protege frente a usos directos del transformador al margen del detector.

    \item \textbf{Clonación de las condiciones.}
    Se clonan \lstinline|A| y \lstinline|B| antes de reutilizarlas en la
    condición combinada. El \emph{aliasing} es la situación en que dos partes del
    árbol comparten el mismo nodo: como en JavaParser cada nodo tiene un único
    nodo padre, insertar directamente la condición original en el nuevo
    \lstinline|if| la desligaría de su ubicación previa y dejaría el AST en un
    estado inconsistente. Clonar produce copias independientes y evita ese efecto.

    \item \textbf{Parentizado de subexpresiones con menor precedencia.}
    Cuando alguna condición es una expresión \lstinline=||=, se encapsula mediante \lstinline|EnclosedExpr| para preservar correctamente la precedencia frente a \lstinline|&&|.

    \item \textbf{Construcción de la condición combinada.}
    Se crea un nuevo \lstinline|BinaryExpr| con operador \lstinline|AND|, manteniendo el orden original de evaluación como \lstinline|outer && inner|.

    \item \textbf{Sustitución del nodo externo.}
    La condición del \lstinline|if| externo se reemplaza por la nueva condición combinada y el cuerpo pasa a ser el del \lstinline|if| interno, previamente clonado y normalizado a \lstinline|BlockStmt|. Por último se invoca \lstinline|removeElseStmt()| como salvaguarda adicional (\emph{no-op} si, como garantiza la guarda previa, ya no hay \lstinline|else|).
\end{enumerate}
\subsection{Política iterativa una por una}
\label{sec:one-at-a-time}

Las múltiples oportunidades dentro de un mismo método se procesan \textbf{de una en una}: tras cada transformación se vuelve a ejecutar la detección sobre el AST modificado, hasta que no queden más oportunidades o se alcance un límite de seguridad de 10 pasadas (\lstinline|MAX_PASSES| en \texttt{BatchRunner}), un tope que garantiza la terminación del bucle de re-detección aunque, por un error inesperado, una pasada no llegara a reducir el número de oportunidades. Esta política se adoptó al observar que aplicar todas las oportunidades en un único pase podía producir sobreconteo cuando los nodos afectados estaban anidados unos dentro de otros.

En concreto, esta decisión garantiza que:

\begin{enumerate}
    \item \lstinline|opportunitiesApplied| refleje únicamente transformaciones realmente aplicadas;
    \item la métrica $\Delta (\Delta = after - before)$ calculada antes y después de la refactorización sea exacta para cada caso; y
    \item la trazabilidad del proceso iterativo quede registrada mediante \lstinline|totalPasses|.
\end{enumerate}

En conjunto, el diseño sacrifica cobertura a cambio de garantías: los casos con
\lstinline|else| o con sentencias adicionales en el bloque externo quedan fuera
de alcance, pues no son combinables sin alterar el comportamiento. Esta
restricción, puramente estructural, permite que el pipeline determinista actúe
como \textbf{baseline verificable} frente a los LLMs en RQ3.
